% !PRINTED 25
\weberpage{25}

%<<CONT
folgende Glied $Q_1$ nicht Normaltheiler von $P$ ist. Dann giebt
es unter den nach $P$ conjugirten Gruppen $a^{-1} Q_1 a$ mehrere verschiedene,
die wir mit

\weq{1}{Q_1,\ Q_1',\ Q_1'',\ Q_1''' \cdots}

bezeichnen wollen. \gesperrt{Der grösste gemeinschaftliche Theiler
$R$ aller dieser Gruppen ist wieder Normaltheiler von $P$.}

Nun sind aber alle die Gruppen (1) grösste Normaltheiler
von $Q$, wie sich nach dem Schlusssatze des \S.\,3 daraus ergiebt,
dass $Q, Q_1$ mit $a^{-1} Q a, a^{-1} Q_1 a$ isomorph sind, während $Q$ als
Normaltheiler von $P$ mit $a^{-1} Q a$ identisch ist. Wir können also
in der Compositionsreihe auf $Q$ jede der Gruppen (1) folgen
lassen, die alle denselben Index $\nu$ haben.

Ist nun $Q_2$ der grösste gemeinschaftliche Theiler von $Q_1, Q_1'$
(also von $Q_1$ verschieden), so ist $Q_2$ nach dem Satze \S.\,6, 4. ein
grösster Normaltheiler von $Q_1$ und $Q_1'$, dessen Index gleichfalls
$\nu$ ist. Also können wir die Compositionsreihe von $Q$ an auf die
folgenden beiden Arten fortsetzen:

\weq{2}{
\begin{array}{ccc@{\qquad}ccc}
Q, & Q_1, & Q_2, & Q & Q_1' & Q_2 \\
& \nu & \nu & & \nu & \nu
\end{array}
}

Das hierin ausgedrückte Gesetz wollen wir nun verallgemeinern
und durch vollständige Induction beweisen.

Wenn $Q_2$ noch nicht gleich $R$ ist, so fügen wir zu $Q_1, Q_1'$
eine dritte Gruppe $Q_1''$ aus der Reihe (1) hinzu, so dass der
Durchschnitt $Q_3$ von $Q_1, Q_1', Q_1''$ ein echter Theiler von $Q_2$ wird,
und fahren so fort, bis wir zum grössten gemeinschaftlichen
Theiler $R$ aller Gruppen (1) gelangt sind.

Es sei also:

\weq{3}{
\begin{array}{l}
Q_2 \text{ der Durchschnitt von } Q_1,\ Q_1' \\
Q_3 \text{ \ ,,\qquad ,,\qquad ,,\ } Q_1,\ Q_1',\ Q_1'' \\
\cdot\ \cdot\ \cdot\ \cdot\ \cdot\ \cdot\ \cdot\ \cdot\ \cdot\ \cdot\ \cdot\ \cdot\ \cdot \\
Q_r \text{ \ ,,\qquad ,,\qquad ,,\ } Q_1,\ Q_1',\ Q_1'' \cdots Q_1^{(r-1)}
\end{array}
}

und die Anordnung der Gruppen $Q_1, Q_1' \cdots Q_1^{(r-1)}$ sei so gewählt,
dass von den Gruppen $Q_1, Q_2 \cdots Q_r$ jede ein echter
Theiler der vorangehenden ist. Diese Anordnung lässt sich fortsetzen,
so lange $Q_r$ nicht gleich $R$ ist.

Wir wollen beweisen, dass wir die Compositionsreihe von $P$
so wählen können, dass darin

\weq{4}{
\begin{array}{cccc}
Q, & Q_1, & Q_2 & \cdots\ Q_r\\
& \nu' & \nu & \nu % [sic: the source prints \nu' although the text says all indices are \nu.]
\end{array}
}
%CONT>>
